\subsection{Products of arbitrary finite
blocks}\label{proof:positive_structure:products-of-arbitrary-finite-blocks}

Partition the coordinates into arbitrary disjoint blocks of size at most
\(B\). Let

\[
\mu=\bigotimes_b\pi_b,
\qquad\pi_b(x_b)\ge a>0
\quad\text{for every block atom.}
\tag{BP}
\]

The block laws may differ and depend on the total dimension; \(B,a\) are
uniform.

\textbf{Theorem P3 (bounded blocks).} There is \(c_{B,a}>0\) such that

\[
\boxed{\max_i I_i^\mu(f)
\ge c_{B,a}\frac{v^{5/2}}{d^6\log^6(2d/v)}
\ge c'_{B,a}\frac{v^{7/2}}{d^7}.}
\tag{P3}
\]

For \(v=0\), use the polynomial bound with right side zero. Arbitrary
real atom probabilities are permitted: no finite-precision assumption is
imposed.

\subsubsection{Step 1: simultaneous dyadic approximation avoids
dimension
error}\label{proof:positive_structure:step-1-simultaneous-dyadic-approximation-avoids-dimension-error}

Approximate each block law by a law \(\pi'_b\) with denominator \(2^k\),
total mass one, and atom errors at most \(2^{-k}\). Round masses down
and distribute the missing units. If \(2^{-k}\le a/2\), every new atom
is at least \(a/2\).

Put \(\pi_{b,t}=\pi_b+t(\pi'_b-\pi_b)\). These are probability laws whenever
\(|t|\le R=a2^{k-1}\): every atom remains at least \(a/2\), and the
masses sum to one. For any cube-bounded multilinear degree-\(r\)
function \(h\),

\[
q_h(t)=\mathbb E_{\otimes_b\pi_{b,t}}h
\]

is a polynomial of degree at most \(r\): a monomial touches at most
\(r\) blocks. It is bounded by one on \([-R,R]\). When \(R\ge 1\),
Markov's inequality gives

\[
|\mathbb E_{\mu'}h-\mathbb{E}_\mu h|
\le r^2/R=2r^2 2^{-k}/a.
\tag{3.1}
\]

There is no dependence on the number of blocks. Applying this to
\(f,f^2,(D_i f)^2\), with multilinearization where necessary, yields

\[
|\operatorname{Var}_{\mu'}f-v|\le12d^2 2^{-k}/a,
\quad |I_i^{\mu'}(f)-I_i^\mu(f)|\le8d^2 2^{-k}/a.
\tag{3.2}
\]

\subsubsection{Step 2: exact sampling of approximate
blocks}\label{proof:positive_structure:step-2-exact-sampling-of-approximate-blocks}

Sample each \(\pi'_b\) with exactly \(k\) fair bits, assigning the
\(2^k\) source strings with the prescribed multiplicities. Every output
coordinate has degree at most \(k\). Hence \(F=f\circ T\) has degree at
most \(kd\) and variance
\(\operatorname{\operatorname{Var}}_{\mu\prime} f\).

A source-bit flip changes only its own block, at most \(B\) outputs.
Every intermediate output law is at most \((2/a)\mu '\): outside the
block the law is unchanged and independent, and every block atom has
mass at least \(a/2\). The telescoping proof of P1 gives

\[
\max_i I_i^{\mu'}(f)
\ge\frac{c_0a}{2B^2}
\frac{\operatorname{Var}_{\mu'}(f)^{5/2}}{k^6d^6}.
\tag{3.3}
\]

This direct full-block comparison is stronger than applying a
single-site floor for each changed coordinate.

\subsubsection{Step 3: choose precision using degree and
variance}\label{proof:positive_structure:step-3-choose-precision-using-degree-and-variance}

Choose \(k=\lceil A_{B,a}\log_2(2d/v)\rceil\), where \(A_{B,a}\) is
sufficiently large. Then \(R\ge 1\), the variance error in (3.2) is at
most \(v/2\), and its influence error is at most half of (3.3) with
variance replaced by \(v/2\).

To justify a uniform choice, put \(x=2d/v\ge 2\). Since
\(d^8/v^{5/2}\le x^8\), the needed error comparison reduces to
\(2^{-k}k^6\le c_{B,a}x^{-8}\). With \(k\le A_{B,a}\log _2 x+1\),
sufficiently large fixed \(A_{B,a}\) makes this true for all \(x\ge 2\).
Exponential decay in \(k\) dominates its sixth power. No choice depends
separately on dimension.

Equations (3.2)--(3.3) prove the first bound in P3. The universal bound
\(\log ^6(2d/v)\le C d/v\) gives the second. \qed

\subsubsection{Interaction interpretation and a density-method
counterexample}\label{proof:positive_structure:interaction-interpretation-and-a-density-method-counterexample}

If a classical Hamiltonian decomposes over components of size at most
\(B\), each with energy oscillation at most \(W\), its Gibbs law
satisfies (BP) with \(a\ge 2^{-B}e^{-\beta  W}\). Thus fixed finite
inverse temperature, bounded component size, and bounded component
energy oscillation imply AA. There is no universal finite-temperature
failure threshold in this class; constants may deteriorate as
temperature approaches zero.

For \(m\) independent Ising edges with coupling \(K=\beta  J \ne 0\),
relative density to uniform is a product of \(e^{Kxy}/\cosh  K\). Its
minimum is \((e^{-|K|}/\cosh  K)^m\), tending exponentially to zero. For
every fixed \(p>1\),

\[
\left\|\frac{d\mu}{dU}\right\|_p
=\left(\frac{\cosh(pK)}{\cosh(K)^p}\right)^{m/p},
\]

which grows exponentially. P3 still proves AA at every fixed finite
coupling. Independent biased spins are the block-size-one special case.

\subsubsection{An alternative finite-alphabet
reduction}\label{proof:positive_structure:an-alternative-finite-alphabet-reduction}

A \(K\le 2^B\)-atom block can also be sampled by \(K-1\) independent
stick-breaking Bernoulli variables, with stage-\(j\) probability
\(\pi_j/(\pi_j+\cdots+\pi_K)\). If all atoms are at least \(a\), these
probabilities lie in \([a,1-a]\). Output coordinates have degree at most
\(K-1\). The block-path argument then transfers AA from bounded-bias
products.
\href{https://arxiv.org/html/2609.35327v1}{Girish--Mittal--Nehoran--Raz,
Remark 4.5} explicitly permits coordinate-dependent biases bounded away
from zero and one. The proof above avoids this additional imported
theorem.
